\newcolumntype{L}[1]{>{\raggedright\arraybackslash\hsize=#1\hsize}X}
\newcolumntype{C}[1]{>{\centering\arraybackslash\hsize=#1\hsize}X}
\begin{landscape}
\begin{table}
\centering
\begin{threeparttable}
\scriptsize
\setlength{\tabcolsep}{4pt}
\caption{Example table showcasing the use of landscape, threeparttable, tabularx, and custom columns.}
\label{tab:example_overview}

\begin{tabularx}{\textwidth}{@{} L{1.20} L{1.30} L{0.60} L{1.55} L{0.85} @{}}
\toprule
\textbf{Item} &
\textbf{Category} &
\textbf{Type} &
\textbf{Description} &
\textbf{Status} \\
\midrule

Item A &
Category 1 &
Type X &
Description of Item A &
Active\tnote{*} \\

Item B &
Category 2 &
Type Y &
Description of Item B &
Inactive \\

\bottomrule
\end{tabularx}

\begin{tablenotes}
  \item[*] Example footnote explaining the active status.
\end{tablenotes}

\end{threeparttable}
\end{table}
\end{landscape}
